\documentclass[11pt]{article}
\usepackage{mathblatt}
% Probe 08.10.2026: „Die lange Seite“ in drei Durchgängen, Durchgang 1 (Einstieg), von Hand gesetzt
% nach aufgabenbank bau/bauregeln.md (Stand 08.10.2026) und mathe-nachhilfe plan.md § 5 „Richtung Durchgänge“.
% Einstieg: mit Vorstufen, Hilfen aufgedeckt, lange Päckchen, glatte Zahlen. Satz aus paket-pythagoras-2.
% Gemeinsame Präambel der drei Durchgänge, übernommen aus paket-pythagoras-2/2-blatt.tex (Probe 08.10.2026).
\makeatletter
\setlength{\mbpfbe}{0mm}% keine Punkte (Bauregel 6.4)
% eingerückt (Bauregel 4.7, 6.6): \gEin Einzug, \gSchrift eine Stufe kleiner
\newlength{\gEin}\setlength{\gEin}{0pt}
\let\gSchrift\relax
\renewcommand{\pfkreuzzeile}[1]{\par\noindent\hangindent3.2em\hangafter1\hspace*{1.6em}$\square$\ #1\par}
\newcommand{\gkreuz}[1]{$\square$\ #1\hspace{2em}}
% Bauregel 6.7: links, was man ansieht, rechts, was man tut; Spaltengrenze überall an derselben Stelle
\newsavebox{\gbox}\newlength{\gLinks}
\newcommand{\gzwei}[2]{\par\smallskip\noindent\sbox{\gbox}{#1}%
  \setlength{\gLinks}{\dimexpr0.5\textwidth-\gEin-\mbpfnr\relax}%
  \ifdim\wd\gbox>\dimexpr\gLinks-4mm\relax
    \usebox{\gbox}\par\smallskip\noindent\begin{minipage}{\linewidth}#2\end{minipage}%
  \else
    \begin{minipage}[t]{\gLinks}\vspace{0pt}\usebox{\gbox}\end{minipage}%
    \begin{minipage}[t]{\dimexpr\linewidth-\gLinks\relax}\vspace{0pt}\raggedright #2\end{minipage}%
  \fi\par}
\RenewDocumentEnvironment{pfaufg}{m m m}{%
  \begin{lrbox}{\mbaufgabebox}\begin{minipage}[t]{\linewidth}%
  \gSchrift\noindent\hspace*{\gEin}\mbpfrandmarke{#1}%
  \begin{minipage}[t]{\mbpfnr}\strut\textbf{#2}\end{minipage}%
  \begin{minipage}[t]{\dimexpr\linewidth-\gEin-\mbpfnr\relax}\raggedright\strut\ignorespaces}%
 {\end{minipage}%
  \end{minipage}\end{lrbox}%
  \setlength{\mbaufgabehoehe}{\dimexpr\ht\mbaufgabebox+\dp\mbaufgabebox\relax}%
  \par\addvspace{7pt}\Needspace*{\mbaufgabehoehe}%
  \noindent\usebox{\mbaufgabebox}\par\addvspace{7pt}}
\makeatother
% Rechenraum als Karo (Bauregel 6.9): n Rechenschritte = n mal 10 mm
\newcommand{\karo}[1]{\par\smallskip\noindent\begin{tikzpicture}
  \clip (0,0) rectangle ({\linewidth-1mm},{#1*10mm});
  \draw[black!16,line width=0.3pt,step=5mm] (0,0) grid ({\linewidth},{#1*10mm});
  \end{tikzpicture}\par}
\newcommand{\kk}[1]{$\square$\,#1\hspace{0.9em}}
\newcommand{\linie}[1]{\,{\color{black!45}\rule[-1pt]{#1}{0.4pt}}\,}
% Zeile am Ende jedes Durchgangs: Originale klein grau, darüber; dann Vorgänger/Nachfolger (Bauregel 3.4, 6.5)
\newcommand{\blattende}[2]{\par\vfill\noindent\begin{minipage}{\linewidth}{\scriptsize\color{mbgrau}Originale: #1\par}\smallskip
{\footnotesize #2\par}\end{minipage}}
\newcommand{\pkt}{\textperiodcentered{}}

\newcommand{\kennung}{H6Q}
% Päckchenzeile mit Lücke für c² (Hilfe aufgedeckt)
\newcommand{\pz}[3]{\noindent\makebox[47mm][l]{#1}$c^2 =$\linie{42mm}\hspace{5mm}$c #2$\linie{18mm}#3\par\medskip}

\begin{document}
\pfheftstil
\fancyfoot[C]{\parbox[t]{\textwidth}{\mbfhbox\par\vspace{1.5mm}{\footnotesize\color{mbgrau}{\scriptsize \kennung}\hfill\thepage}}}
\pfheftkopf{Die lange Seite \pkt{} Hypotenuse}{P10 \pkt{} 1}

% 1 Vorstufe: Idee sichtbar machen (Bauregel 4.1; Kette 3, Zahlen 3-4-5)
\begin{pfaufg}{}{1.}{}
Das Dreieck hat einen rechten Winkel. Über jeder Seite steht ein Quadrat. Die schräge Seite ist 5 Kästchen lang.
\gzwei{\begin{tikzpicture}[scale=0.40,line width=0.7pt]
\draw[mbgitter,line width=0.3pt] (-4,-3) grid (7,7);
\draw (0,0) -- (3,0) -- (3,-3) -- (0,-3) -- cycle;
\draw (0,0) -- (-4,0) -- (-4,4) -- (0,4) -- cycle;
\draw (0,4) -- (3,0) -- (7,3) -- (4,7) -- cycle;
\draw[line width=1.4pt] (0,0) -- (3,0) -- (0,4) -- cycle;
\mbrwmarke{(0,0)}{(3,0)}{(0,4)}
\end{tikzpicture}}{%
\textbf{a)}\ Zähle die Kästchen in jedem Quadrat.\par\smallskip
\hspace*{1.4em}klein:\linie{12mm}und\linie{12mm}\quad groß:\linie{12mm}\par\bigskip
\textbf{b)}\ Kreuze an, was stimmt.
\pfkreuzzeile{Die zwei kurzen Seiten sind zusammen so lang wie die lange Seite.}
\pfkreuzzeile{Die zwei kleinen Quadrate haben zusammen so viele Kästchen wie das große.}}
\fusshilfe{\mbox{1:~a) 9, 16, 25; b) die Quadrate}}
\end{pfaufg}

% 2 Vorstufe: Erkennen in vier Lagen; d) ohne rechten Winkel (Kette 1 s0 v3: stumpfer Winkel)
\begin{pfaufg}{}{2.}{}
Die lange Seite gegenüber dem rechten Winkel heißt \textbf{Hypotenuse}. Die zwei kurzen Seiten am rechten Winkel heißen \textbf{Katheten}. Kreuze in jedem Dreieck die Hypotenuse an. Hat ein Dreieck keinen rechten Winkel, kreuze „keine“ an.\par\medskip\noindent
\begin{minipage}[b]{0.245\linewidth}\centering
\begin{tikzpicture}[line width=0.7pt,scale=0.85]\coordinate (P1) at (0,0);\coordinate (P2) at (2.5,0);\coordinate (P3) at (0.9,1.2);\draw (P1) -- (P2) -- (P3) -- cycle;\mbrwmarke{(P3)}{(P1)}{(P2)}
\node[font=\small,anchor=north] at (1.25,-0.08) {$r$};\node[font=\small,anchor=south east] at (0.40,0.62) {$s$};\node[font=\small,anchor=south west] at (1.75,0.62) {$t$};\end{tikzpicture}\par\smallskip
{\small a)\ \kk{$r$}\kk{$s$}\kk{$t$}\par\kk{keine}}
\end{minipage}\hfill
\begin{minipage}[b]{0.245\linewidth}\centering
\begin{tikzpicture}[line width=0.7pt,scale=0.85]\coordinate (P1) at (0,0);\coordinate (P2) at (2.2,0);\coordinate (P3) at (2.2,1.5);\draw (P1) -- (P2) -- (P3) -- cycle;\mbrwmarke{(P2)}{(P1)}{(P3)}
\node[font=\small,anchor=north] at (1.1,-0.08) {$a$};\node[font=\small,anchor=west] at (2.28,0.75) {$c$};\node[font=\small,anchor=south east] at (1.05,0.80) {$b$};\end{tikzpicture}\par\smallskip
{\small b)\ \kk{$a$}\kk{$b$}\kk{$c$}\par\kk{keine}}
\end{minipage}\hfill
\begin{minipage}[b]{0.245\linewidth}\centering
\begin{tikzpicture}[line width=0.7pt,scale=0.85]\coordinate (P1) at (0,1.25);\coordinate (P2) at (1.8,0);\coordinate (P3) at (2.8,1.44);\draw (P1) -- (P2) -- (P3) -- cycle;\mbrwmarke{(P2)}{(P1)}{(P3)}
\node[font=\small,anchor=north east] at (0.85,0.56) {$f$};\node[font=\small,anchor=north west] at (2.37,0.67) {$g$};\node[font=\small,anchor=south] at (1.40,1.42) {$e$};\end{tikzpicture}\par\smallskip
{\small c)\ \kk{$e$}\kk{$f$}\kk{$g$}\par\kk{keine}}
\end{minipage}\hfill
\begin{minipage}[b]{0.245\linewidth}\centering
\begin{tikzpicture}[line width=0.7pt,scale=0.85]\coordinate (P1) at (0.4,0);\coordinate (P2) at (2.9,0);\coordinate (P3) at (0,1.1);\draw (P1) -- (P2) -- (P3) -- cycle;
\node[font=\small,anchor=north] at (1.65,-0.08) {$u$};\node[font=\small,anchor=east] at (0.15,0.5) {$v$};\node[font=\small,anchor=south west] at (1.5,0.58) {$w$};\end{tikzpicture}\par\smallskip
{\small d)\ \kk{$u$}\kk{$v$}\kk{$w$}\par\kk{keine}}
\end{minipage}
\fusshilfe{\mbox{2:~a) $r$; b) $b$; c) $e$; d) keine}}
\end{pfaufg}

% 3 Satz in Worten, P10-Form (2022 · 1 · g)
\begin{pfaufg}{P10 ’22}{3.}{}
Kreuze die Aussage an, die für jedes rechtwinklige Dreieck stimmt.
\pfkreuzzeile{Dem rechten Winkel liegt eine Kathete gegenüber.}
\pfkreuzzeile{$c = a + b$}
\pfkreuzzeile{Die Summe der Flächeninhalte der Kathetenquadrate ist gleich dem Flächeninhalt des Hypotenusenquadrates.}
\fusshilfe{\mbox{3:~die dritte Aussage}}
\end{pfaufg}

% 4 Gleichung aufstellen mit Hilfe (Bauregel 4.5): z ist die Hypotenuse (2017 · 1 · d)
\begin{pfaufg}{P10 ’17}{4.}{}
Kreuze die Gleichung an, die für dieses Dreieck gilt.\par
Suche zuerst die Hypotenuse – sie steht allein auf einer Seite.
\gzwei{\begin{tikzpicture}[line width=0.7pt]\coordinate (P1) at (0,0);\coordinate (P2) at (3.2,0);\coordinate (P3) at (0,2.2);\draw (P1) -- (P2) -- (P3) -- cycle;\mbrwmarke{(P1)}{(P2)}{(P3)}
\node[font=\small,anchor=north] at (1.6,-0.08) {$y$};\node[font=\small,anchor=east] at (-0.08,1.1) {$x$};\node[font=\small,anchor=south west] at (1.62,1.12) {$z$};\end{tikzpicture}}{%
\pfkreuzzeile{$x^2 = y^2 + z^2$\hspace{3em}$\square$\ $z^2 = x^2 - y^2$}
\pfkreuzzeile{$z^2 \cdot y^2 = x^2$\hspace{3.15em}$\square$\ $z^2 = x^2 + y^2$}}
\fusshilfe{\mbox{4:~$z^2 = x^2 + y^2$}}
\end{pfaufg}

% 5 Formel mit Wurzel, ohne Hilfe (2021 · 1 · h)
\begin{pfaufg}{P10 ’21}{5.}{}
Kreuze die Formel an, mit der man $z$ berechnet.
\gzwei{\begin{tikzpicture}[line width=0.7pt]\coordinate (P1) at (0,0);\coordinate (P2) at (3.2,0);\coordinate (P3) at (3.2,2.2);\draw (P1) -- (P2) -- (P3) -- cycle;\mbrwmarke{(P2)}{(P1)}{(P3)}\node[font=\small,anchor=west] at (3.28,1.1) {$y$};\node[font=\small,anchor=south east] at (1.555,1.166) {$z$};\node[font=\small,anchor=north] at (1.6,-0.08) {$x$};\end{tikzpicture}}{%
\pfkreuzzeile{$z = \sqrt{x^2 - y^2}$\hspace{2.4em}$\square$\ $z = \sqrt{x^2 + y^2}$}
\pfkreuzzeile{$z = \sqrt{y^2 - x^2}$\hspace{2.4em}$\square$\ $z = y^2 + x^2$}}
\fusshilfe{\mbox{5:~$z = \sqrt{x^2 + y^2}$}}
\end{pfaufg}

% 6 langes Päckchen (Bauregel 4.7): Kette 2 s1 v1, s1 v3, s4 v1, s5 v1, s3 v1; a) eigene Kopfrechenzahl
% Hilfe aufgedeckt: die Lücke zeigt die Zwischengröße c²
\begin{pfaufg}{}{6.}{}
Berechne die Hypotenuse $c$ aus den beiden Katheten. Runde in f) auf eine Stelle nach dem Komma.\par\medskip
\pz{a)\ $6$ cm und $8$ cm}{=}{cm}
\pz{b)\ $36$ cm und $15$ cm}{=}{cm}
\pz{c)\ $36$ cm und $48$ cm}{=}{cm}
\pz{d)\ $2{,}5$ m und $6$ m}{=}{m}
\pz{e)\ $1{,}2$ m und $50$ cm}{=}{cm}
\pz{f)\ $4$ cm und $9$ cm}{\approx}{cm}
\fusshilfe{\mbox{6:~a) 10 cm; b) 39 cm; c) 60 cm; d) 6,5 m; e) 130 cm; f) $\sqrt{97} \approx 9{,}8$ cm}}
\end{pfaufg}

% 7 Dreieck in anderer Lage (Kette 2 s2 v1): die Hypotenuse liegt unten waagerecht; Hilfe a) deckt die Rolle auf
\begin{pfaufg}{}{7.}{}
Im Dreieck ist eine Seite mit ? markiert.
\gzwei{\begin{tikzpicture}[line width=0.7pt,scale=0.8]\coordinate (A) at (0,0);\coordinate (B) at (5,0);\coordinate (C) at (1.4,2.24);\draw (A) -- (B) -- (C) -- cycle;\mbrwmarke{(C)}{(A)}{(B)}
\node[font=\small,anchor=north] at (2.5,-0.08) {?};\node[font=\small,anchor=south east] at (0.66,1.12) {$28$ cm};\node[font=\small,anchor=south west] at (3.25,1.15) {$45$ cm};\end{tikzpicture}}{%
\textbf{a)}\ Die Seite ? ist \kk{eine Kathete}\kk{die Hypotenuse}\par\medskip
\textbf{b)}\ Berechne die Länge der Seite ?.
\karo{3}\par\smallskip\hfill ? $=$\linie{18mm}cm}
\fusshilfe{\mbox{7:~a) die Hypotenuse; b) 53 cm}}
\end{pfaufg}

% 8 Diagonale im Rechteck (Kette 2 s9 v1): die Skizze zeigt das Dreieck
\begin{pfaufg}{}{8.}{}
Wie lang ist die Diagonale $d$ des Bilderrahmens?
\gzwei{\begin{tikzpicture}[line width=0.7pt]\draw (0,0) rectangle (3.6,1.92);\draw (0.25,0.25) rectangle (3.35,1.67);\draw[line width=1.2pt] (0,0) -- (3.6,1.92);
\node[font=\small,anchor=north] at (1.8,-0.08) {$45$ cm};\node[font=\small,anchor=west] at (3.68,0.96) {$24$ cm};\node[font=\small,anchor=south east] at (1.65,0.95) {$d$};\end{tikzpicture}}{%
\karo{3}\par\smallskip\hfill $d =$\linie{18mm}cm}
\fusshilfe{\mbox{8:~$d = 51$ cm}}
\end{pfaufg}

% 9 Sache mit Skizze (Kette 2 s10 v3): Seil am Mast; Maße nur in der Skizze (Bauregel 6.10)
\begin{pfaufg}{}{9.}{}
Ein Seil hält einen Mast. Wie lang ist das Seil? Runde auf eine Stelle nach dem Komma.
\gzwei{\begin{tikzpicture}[line width=0.7pt,scale=0.45]
\draw (-0.6,0) -- (5.2,0);
\draw[line width=2pt] (4,0) -- (4,8.6);
\draw (0,0) -- (4,8);\fill (4,8) circle (2pt);
\mbrwmarke{(4,0)}{(0,0)}{(4,8)}
\node[font=\small,anchor=north] at (2,-0.1) {$3$ m};
\node[font=\small,anchor=west] at (4.2,4) {$8$ m};
\node[font=\small,anchor=south east] at (1.9,4.1) {Seil};
\end{tikzpicture}}{%
\karo{3}\par\smallskip\hfill Seil $\approx$\linie{18mm}m}
\fusshilfe{\mbox{9:~$\sqrt{73} \approx 8{,}5$ m}}
\end{pfaufg}

% 10 Überstand (Kette 2 s11 v1); Hilfe deckt die Zwischengröße auf (Bauregel 4.4)
\begin{pfaufg}{}{10.}{}
Eine Leiter lehnt an einer Mauer. Oben ragt sie über die Mauerkante hinaus. Wie lang ist die ganze Leiter?\par
Berechne zuerst das schräge Stück bis zur Mauerkante.
\gzwei{\begin{tikzpicture}[line width=0.7pt,scale=0.55]
\fill[black!10] (2.4,0) rectangle (3.0,4);\draw (2.4,0) -- (2.4,4) -- (3.0,4);
\draw (-0.6,0) -- (3.6,0);
\draw[line width=1.6pt] (0,0) -- (2.985,4.976);
\mbrwmarke{(2.4,0)}{(0,0)}{(2.4,4)}
\draw[|-|,line width=0.4pt] (0,-0.45) -- (2.4,-0.45) node[midway,below,font=\small] {$0{,}9$ m};
\node[font=\small,anchor=west] at (3.05,2.0) {$4$ m};
\node[font=\small,anchor=east] at (2.55,4.6) {$1$ m};
\node[font=\scriptsize,mbgrau,anchor=north west] at (-0.6,-1.3) {nicht maßstabsgerecht};
\end{tikzpicture}}{%
\karo{4}\par\smallskip\hfill Leiter:\linie{18mm}m}
\fusshilfe{\mbox{10:~$4{,}1 + 1 = 5{,}1$ m}}
\end{pfaufg}

% 11 Fehler finden (Kette 5 s1 v1): Wurzel vergessen
\begin{pfaufg}{}{11.}{}
Lara soll die Hypotenuse $c$ berechnen. Die Katheten sind $33$~cm und $56$~cm lang. Finde ihren Fehler und rechne richtig.
\gzwei{\hspace*{4mm}$\begin{aligned} c^2 &= 33^2 + 56^2\\ c^2 &= 4\,225\\ c &= 4\,225 \text{ cm}\end{aligned}$}{%
Fehler:\linie{50mm}\par\bigskip
\karo{2}\par\smallskip\hfill $c =$\linie{18mm}cm}
\fusshilfe{\mbox{11:~Wurzel vergessen; $c = 65$ cm}}
\end{pfaufg}

% 12 Prüfungshöhe (Kette 2 s15 v1, eigene Variante zu 2020 · 7 · a): Antwortzeile in Form des Ergebnisses
\begin{pfaufg}{}{12.}{}
Ein Grundstück $ABC$ hat bei $C$ einen rechten Winkel. Berechne die Länge der Seite $\overline{AB}$. Runde auf eine Stelle nach dem Komma.
\gzwei{\begin{tikzpicture}[line width=0.7pt]\coordinate (C) at (0,0);\coordinate (B) at (1.26,0);\coordinate (A) at (0,3.51);\draw (A) -- (B) -- (C) -- cycle;\mbrwmarke{(C)}{(B)}{(A)}
\node[font=\small,anchor=north east] at (C) {$C$};\node[font=\small,anchor=north west] at (B) {$B$};\node[font=\small,anchor=south] at (A) {$A$};
\node[font=\small,anchor=north] at (0.63,-0.08) {$14$ m};\node[font=\small,anchor=east] at (-0.08,1.75) {$39$ m};\end{tikzpicture}}{%
\karo{3}\par\smallskip\hfill $\overline{AB}^2 =$\linie{22mm}\quad $\overline{AB} \approx$\linie{18mm}m}
\fusshilfe{\mbox{12:~$\overline{AB}^2 = 1\,717$, $\overline{AB} \approx 41{,}4$ m}}
\end{pfaufg}

\blattende{2017 \pkt{} 1 \pkt{} d \quad 2021 \pkt{} 1 \pkt{} h \quad 2022 \pkt{} 1 \pkt{} g}{Vorher: Quadrat und Wurzel\quad\textbullet\quad Weiter: Die lange Seite \pkt{} P10 \pkt{} 2}
\end{document}
